\documentclass[11pt, a4paper]{article}
\usepackage[left=2cm,right=2cm,top=2cm,bottom=2cm]{geometry}
\usepackage{graphicx} % Required for inserting images
\usepackage{amsmath}
\usepackage[dvipsnames]{xcolor}

\usepackage{makecell}
\usepackage{tabularx}
% center shit vertically
\def\tabularxcolumn#1{m{#1}}

% \usepackage[a4paper, total={6in, 8in}]{geometry}

\title{Test Plan for Shadow Operative}
\date{\today}

\author{
    Ben Magowan \\
    \texttt{vcwx72}
    \and
    Danial Maqbool \\
    \texttt{gkzv63}
    \and
    Will Morgan \\
    \texttt{sghp78}
    \and
    Edmond Vajda \\
    \texttt{njql73} \\
    \and
    Sami Lodi \\
    \texttt{hqcb32}
    \and
    Heria Chen \\
    \texttt{qzxl76}
    \and
    Toby Davis \\
    \texttt{cltz62}
}

\begin{document}

\maketitle

\begin{center}
    \text{Group Number 05} \\
    \text{Version 1.0.1}
\end{center}

\newpage

\tableofcontents

\newpage

\section{Introduction and Test Overview}

\subsection{Introduction}

\noindent Shadow Operative integrates IBM SkillsBuild content into a video game, creating an innovative platform that makes learning about AI, Cybersecurity, and Data Analytics more engaging for younger audiences. The game incorporates IBM SkillsBuild badges and leverages AI-driven adaptive learning to personalize content based on each player's performance and progress. This approach blends education with entertainment, fostering a learning experience that is not only effective but also enjoyable. \\

\noindent The integration ensures players can earn recognizable certifications while playing, enhancing their skills in high-demand areas through practical, game-based challenges. As players advance, the game dynamically adjusts difficulty and content, tailoring lessons to their unique needs and encouraging mastery of complex topics. \\

\subsection{Test Overview}

\subsubsection{Unit Testing}

\begin{tabularx}{\linewidth}{|c|X|X|}
    \hline
    \textbf{Test ID} & \textbf{Description} & \textbf{Test Oracle} \\
    \hline
    \textbf{utest\_1} & Character Movement: Test that the character moves correctly in response to input, including sprinting, crouching, and jumping & Expected character movement should match input commands; animations and physics should behave consistently \\
    \hline
    \textbf{utest\_2} & Health: Test that health decreases upon taking damage and increases when healing items are used & Health bar reflects correct changes based on game events \\
    \hline
    \textbf{utest\_3} & Collision Detection: Verify accurate collision between player, environment, and interactive objects & Player and objects cannot pass through solid obstacles; correct response on collision events \\
    \hline
    \textbf{utest\_4} & Enemy AI Behaviour: Test patrol routes, detection logic, and state transitions (e.g., alert, idle) & AI responds correctly to player actions and follows expected behavior patterns \\
    \hline
    \textbf{utest\_5} & Level Load + Transitions: Test smooth transitions between levels and correct loading of level assets & Level data persists, and transitions occur without graphics or gameplay glitches \\
    \hline
    \textbf{utest\_6} & Skill Build Questions: Test that questions display and function correctly during gameplay & Questions render properly, validate player responses, and provide appropriate feedback \\
    \hline
    \textbf{utest\_7} & Level Objectives: Verify completion of primary and secondary objectives & Objectives track progress and provide correct rewards upon completion \\
    \hline
    \textbf{utest\_8} & Sound Effects: Test audio playback for various game events & Correct audio plays at the appropriate volume and timing \\
    \hline
    \textbf{utest\_9} & Environmental Effects: Verify rendering of lighting, shadows, and ambient sounds in levels & Effects render without graphical glitches and match the game environment \\
    \hline
    \textbf{utest\_10} & Dynamic Question Difficulty: Test adaptive difficulty of quiz questions based on player performance & Question difficulty adjusts dynamically based on performance metrics \\
    \hline
    \textbf{utest\_11} & Enemy Health: Test that enemy health decreases correctly when hit by weapons & Enemy health bar updates properly, and enemies die when health reaches zero \\
    \hline
    \textbf{utest\_12} & Weapon Mechanics: Test firing, ammo decrement, and reload functionality & Weapons deal appropriate damage, consume ammo, and reload as expected \\
    \hline
\end{tabularx}

\subsubsection{System Testing}

\begin{tabularx}{\linewidth}{|c|X|X|c|}
    \hline
    \textbf{Test ID} & \textbf{Description} & \textbf{Test Oracle} & \textbf{Nature} \\
    \hline
    \textbf{stest\_1} & Graphics and Audio: Verify that graphics render smoothly, textures load correctly, and audio plays without distortion or lag & Graphics match expected design, and audio aligns with gameplay events without delay & Non-Func \\
    \hline
    \textbf{stest\_2} & UI Navigation (Saving and Loading): Test that the UI allows seamless navigation, and game states save/load correctly & UI responds correctly to input, and saved states are accurately restored & Func \\
    \hline
    \textbf{stest\_3} & Gameplay Flow: Test the overall game experience, including mission progression, level transitions, and story continuity & Gameplay progresses smoothly without logic errors or broken sequences & Func \& N-F \\
    \hline
    \textbf{stest\_4} & Character Interactions: Test interactions between the player and NPCs, objects, and the environment & Characters and objects respond correctly to player interactions with appropriate animations & Func \\
    \hline
    \textbf{stest\_5} & Performance and Stability: Test the game's frame rate, memory usage, and stability under different load conditions & Frame rate remains stable, memory usage is within acceptable limits, and no crashes occur & Non-Func \\
    \hline
    \textbf{stest\_6} & Adaptive AI: Test AI's ability to adjust difficulty and behavior based on player performance & AI dynamically adjusts gameplay difficulty and tactics to challenge the player & Non-Func \\
    \hline
    \textbf{stest\_7} & Educational Content Integration: Test seamless embedding of IBM Skills Build questions into the gameplay & Educational questions appear at appropriate points, and responses impact gameplay as intended & Func \\
    \hline
    \textbf{stest\_8} & Game State Persistence: Verify that in-game events and progress persist across sessions & Game state data is correctly saved and restored across multiple sessions & Func \\
    \hline
    \textbf{stest\_9} & Input Device Compatibility: Test compatibility with different input devices (keyboard, controller) & Game responds consistently to input across various devices & Func \\
    \hline
    \textbf{stest\_10} & Accessibility Features: Test availability and functionality of accessibility features like subtitles and customizable controls & Accessibility options work as intended and improve the user experience & Func \\
    \hline
\end{tabularx}

\newpage

\subsubsection{User Acceptance Testing}

\begin{tabularx}{\linewidth}{|c|X|X|}
    \hline
    \textbf{Test ID} & \textbf{Description} & \textbf{Test Oracle} \\
    \hline
    \textbf{uattest\_1} & Gameplay Experience: Verify that the game provides an engaging, immersive, and challenging experience for players & Players report a fun and balanced experience, with no major gameplay issues \\
    \hline
    \textbf{uattest\_2} & User Interface: Test that the UI is intuitive, responsive, and visually appealing & Players can easily navigate menus, interact with elements, and receive feedback \\
    \hline
    \textbf{uattest\_3} & Performance and Stability: Ensure the game performs smoothly and remains stable under various conditions & Players experience no frame drops, long load times, or crashes \\
    \hline
    \textbf{uattest\_4} & Accessibility: Test that the game is inclusive and accommodates various accessibility needs & Players can use customizable controls, subtitles, and other accessibility features effectively \\
    \hline
    \textbf{uattest\_5} & Quizzes: Test the educational component, ensuring quizzes are engaging, well-integrated, and correctly assessed & Players find the quizzes challenging and educational, with accurate feedback on answers \\
    \hline
    \textbf{uattest\_6} & Chatbot: Test the chatbot's ability to assist users with gameplay tips or educational content & Players receive accurate and helpful responses from the chatbot \\
    \hline
    \textbf{uattest\_7} & Visual and Audio Design: Verify the coherence and appeal of the game's graphics and audio elements & Players find the visual style immersive and the audio enhances the experience \\
    \hline
    \textbf{uattest\_8} & Educational Content Engagement: Test the player's interaction with IBM Skills Build content & Players report an improved understanding of Skills Build concepts through gameplay \\
    \hline
    \textbf{uattest\_9} & User Progress and Rewards: Verify the reward system and progression tracking & Players feel motivated by achievements, and progress is tracked accurately \\
    \hline
    \textbf{uattest\_10} & Difficulty Adjustment: Ensure adaptive AI and difficulty settings enhance player experience & Players report that difficulty levels are well-balanced and adapt to their skill level \\
    \hline
    \textbf{uattest\_11} & Game Narrative and Storyline: Test the clarity and engagement of the storyline & Players find the storyline compelling and easy to follow \\
    \hline
    \textbf{uattest\_12} & Onboarding and Tutorials: Ensure new players understand game mechanics through effective onboarding & Players can learn core mechanics easily and quickly adapt to gameplay \\
    \hline
\end{tabularx}

\newpage

\section{Test Cases}

\subsection{Unit Testing}

\subsubsection{Character Movement}

\begin{tabularx}{\linewidth}{|>{\raggedright\arraybackslash}m{4cm}|X|}
    \hline
    \textbf{Test ID} & \textbf{utest\_1} \\
    \hline
    \textbf{Description of test} & This test case verifies the player's movement controls and response within the game environment. The player should be able to move in all directions (left, right), jump, smoothly and responsively. The controls should be intuitive, and the player's movements should be fluid. \\
    \hline
    \textbf{Related requirement document details} & Movement Mechanics (2.2.8) \\
    \hline
    \textbf{Prerequisites for Test} & \makecell[l]{
        Character controller setup \\
        Basic environment created to navigate, allowing to test horizontal \\ and vertical movement.
    } \\
    \hline
    \textbf{Test Procedure} &
    \textbf{Test character’s horizontal movement}
    \begin{itemize}
        \item Move the character left and right using the input system.
    \end{itemize}
    \textbf{Test character’s vertical movement by jumping midair}
    \begin{itemize}
        \item Press jump while the character is in the air.
    \end{itemize}
    \textbf{Test character’s vertical movement by jumping onto a ledge}
    \begin{itemize}
        \item Press jump when near a ledge to ensure correct landing mechanics.
    \end{itemize}
    \textbf{Test character’s vertical movement by falling off a ledge}
    \begin{itemize}
        \item Walk off the edge of a platform and ensure gravity is applied to the character.
    \end{itemize}
    \textbf{Test character’s movement with simultaneous inputs}
    \begin{itemize}
        \item Press both movement keys (e.g., horizontal and vertical) to ensure combined movement works smoothly.
    \end{itemize}
    \\
    \hline
    \textbf{Test Material Used} & Keyboard, Xbox Controller \\
    \hline
    \textbf{Expected Result} &
    \textbf{Horizontal Movement:}
    \begin{itemize}
        \item The character’s x-coordinate should increase when moving right and decrease when moving left.
        \item Movement should be smooth and responsive.
    \end{itemize}
    \textbf{Vertical Movement (Jumping Mid Air):}
    \begin{itemize}
        \item The character’s y-coordinate should increase when jumping, indicating upward movement.
        \item The jump height should match the intended design.
    \end{itemize}
    \textbf{Vertical Movement (Jumping onto a Ledge):}
    \begin{itemize}
        \item The character’s y-coordinate should decrease to the ledge height, then stop at that height.
        \item The character should land on the ledge, with no clipping or unintended behavior.
    \end{itemize}
    \\
    \hline
\end{tabularx}

\begin{center}
\begin{tabularx}{\linewidth}{|>{\raggedright\arraybackslash}m{4cm}|X|}
    \hline
    \textbf{Expected Result} &
    \textbf{Vertical Movement (Falling off a Ledge):}
    \begin{itemize}
        \item The character’s y-coordinate should decrease when falling off the ledge, simulating gravity.
        \item The fall speed should be consistent with the game’s gravity settings.
    \end{itemize}
    \textbf{Simultaneous Inputs:}
    \begin{itemize}
        \item The character should move both horizontally and vertically in response to input without delay.
        \item No erratic or conflicting movement should occur when both horizontal and vertical inputs are used simultaneously.
    \end{itemize}
    \\
    \hline
    \textbf{Comments} & None \\
    \hline
    \textbf{Created by} & WM \\
    \hline
    \textbf{Test environment(s)} & Windows 11, Mac OS 15.3, Unity 2022.3.50f1 \\
    \hline
\end{tabularx}
\end{center}

\newpage

\noindent \textbf{Results:} This test revealed two regressions: one with Unity's old Input System and another during the upgrade to the new one. We transitioned to the new Unity Input System to streamline the integration of additional input methods. This upgrade is vital for improving accessibility, a key aspect of an educational game, to ensure it accommodates a diverse audience.

\begin{center}
\begin{tabularx}{\linewidth}{|>{\raggedright\arraybackslash}m{4cm}|X|}
    \hline
    \textbf{Test ID} & utest\_1 \\
    \hline
    \textbf{Expected Result} &
    \textbf{Horizontal Movement:}
    \begin{itemize}
        \item The character’s x-coordinate should increase when moving right and decrease when moving left.
        \item Movement should be smooth and responsive.
    \end{itemize}
    \textbf{Vertical Movement (Jumping Mid Air):}
    \begin{itemize}
        \item The character’s y-coordinate should increase when jumping, indicating upward movement.
        \item The jump height should match the intended design.
    \end{itemize}
    \textbf{Vertical Movement (Jumping onto a Ledge):}
    \begin{itemize}
        \item The character’s y-coordinate should decrease to the ledge height, then stop at that height.
        \item The character should land on the ledge, with no clipping or unintended behavior.
    \end{itemize}
    \textbf{Vertical Movement (Falling off a Ledge):}
    \begin{itemize}
        \item The character’s y-coordinate should decrease when falling off the ledge, simulating gravity.
        \item The fall speed should be consistent with the game’s gravity settings.
    \end{itemize}
    \textbf{Simultaneous Inputs:}
    \begin{itemize}
        \item The character should move both horizontally and vertically in response to input without delay.
        \item No erratic or conflicting movement should occur when both horizontal and vertical inputs are used simultaneously.
    \end{itemize}
    \\
    \hline
    \textbf{Regression 1 Result} &
    \textbf{Horizontal Movement:}
    \begin{itemize}
        \item The character’s x-coordinate correctly increases when moving right and decreases when moving left. The movement is smooth and responsive, with no lag or stutter.
    \end{itemize}
    \textbf{Vertical Movement (Jumping Mid Air):}
    \begin{itemize}
        \item The character’s y-coordinate increases when jumping, and the jump height is consistent with the intended design, smoothly reaching the peak before falling.
    \end{itemize}
    \\
    \hline
\end{tabularx}
\end{center}

\begin{center}
\begin{tabularx}{\linewidth}{|>{\raggedright\arraybackslash}m{4cm}|X|}
    \hline
    \textbf{Regression 1 Result} & \textbf{Vertical Movement (Jumping onto a Ledge):}
    \begin{itemize}
        \item The character’s y-coordinate correctly decreases to the ledge height, then stops at that height. The character lands on the ledge without clipping, and there is no unintended behavior (e.g., hovering or falling through the ledge).
    \end{itemize}
    \textbf{Vertical Movement (Falling off a Ledge):}
    \begin{itemize}
        \item The character’s y-coordinate decreases when falling off the ledge, and gravity is applied as expected. The fall speed is consistent with the game’s gravity settings and feels natural.
    \end{itemize}
    \textbf{Simultaneous Inputs:}
    \begin{itemize}
        \item The character moves both horizontally and vertically without delay when both inputs are used. There is no erratic or conflicting movement, and the combined movement feels smooth and fluid.
    \end{itemize}
    \textbf{Unexpected Results}
    \begin{itemize}
        \item No unexpected Results Occurred
    \end{itemize}
    \\
    \hline
    \textbf{Pass/Fail} & Pass \\
    \hline
    \textbf{Regression 2 Result} &
    \textbf{Horizontal Movement:}
    \begin{itemize}
        \item The character’s x-coordinate correctly increases when moving right and decreases when moving left. The movement is smooth and responsive, with no lag or stutter.
    \end{itemize}
    \textbf{Vertical Movement (Jumping Mid Air):}
    \begin{itemize}
        \item The character’s y-coordinate increases when jumping, and the jump height is consistent with the intended design, smoothly reaching the peak before falling.
    \end{itemize}
    \textbf{Vertical Movement (Jumping onto a Ledge):}
    \begin{itemize}
        \item The character’s y-coordinate correctly decreases to the ledge height, then stops at that height. The character lands on the ledge without clipping, and there is no unintended behavior (e.g., hovering or falling through the ledge).
    \end{itemize}
    \textbf{Vertical Movement (Falling off a Ledge):}
    \begin{itemize}
        \item The character’s y-coordinate decreases when falling off the ledge, and gravity is applied as expected. The fall speed is consistent with the game’s gravity settings and feels natural.
    \end{itemize}
    \textbf{Simultaneous Inputs:}
    \begin{itemize}
        \item The character moves both horizontally and vertically without delay when both inputs are used. There is no erratic or conflicting movement, and the combined movement feels smooth and fluid.
    \end{itemize}
    \textbf{Unexpected Results}
    \begin{itemize}
        \item No unexpected Results Occurred
    \end{itemize}
    \\
    \hline
\end{tabularx}
\end{center}

\begin{center}
\begin{tabularx}{\linewidth}{|>{\raggedright\arraybackslash}m{4cm}|X|}
    \hline
    \textbf{Pass/Fail} & Pass \\
    \hline
    \textbf{Severity (If Fail)} & Critical: If this test were to fail, gameplay would be severely impacted, rendering the game unplayable and preventing players from progressing or interacting with game elements. \\
    \hline
    \textbf{Comments} & None \\
    \hline
\end{tabularx}
\end{center}

\subsubsection{Health}

\begin{center}
\begin{tabularx}{\linewidth}{|>{\raggedright\arraybackslash}m{4cm}|X|}
    \hline
    \textbf{Test ID} & utest\_2 \\
    \hline
    \textbf{Description of test} & This test case verifies the health system within the game is affected and responds to stimuli in game. The player should have a number of hearts representing their number of lives remaining. Hearts can be added or taken away depending on certain events in the game. \\
    \hline
    \textbf{Related requirement document details} & Combat Mechanics 2.2.9 \\
    \hline
    \textbf{Prerequisites for test} & Player character available, quizzes and enemies unlocked. \\
    \hline
    \textbf{Test Procedure} &
    \textbf{Test loosing heart via getting question wrong}
    \begin{itemize}
        \item Player character chooses an incorrect answer to a question on screen.
    \end{itemize}
    \textbf{Test loosing heart via getting hit my enemy}
    \begin{itemize}
        \item Player character receives damage from the enemy target.
    \end{itemize}
    \textbf{Test regaining heart via item}
    \begin{itemize}
        \item Player activates an item which restores health.
    \end{itemize}
    \\
    \hline
    \textbf{Test material used} & None \\
    \hline
    \textbf{Expected result} &
    \textbf{Test loosing heart via getting question wrong}
    \begin{itemize}
        \item Player loses 1 heart.
        \item Player’s health bar has a visual indicator of losing a single heart.
    \end{itemize}
    \textbf{Test loosing heart via getting hit by enemy}
    \begin{itemize}
        \item Player loses 1 heart.
        \item Player’s health bar has a visual indicator of losing a single heart.
    \end{itemize}
    \textbf{Test regaining heart via item}
    \begin{itemize}
        \item Player gains 1 heart.
        \item Heart recovered does not go over max hearts.
        \item Player’s health bar has a visual indicator of gaining a heart.
    \end{itemize}
    \\
    \hline
    \textbf{Comments} & None \\
    \hline
    \textbf{Created by} & SL \\
    \hline
    \textbf{Test environment(s)} & Windows 11, Mac OS 15.3, Unity 2022.3.50f1 \\
    \hline
\end{tabularx}
\end{center}

\newpage

\noindent \textbf{Results:}

\begin{center}
\begin{tabularx}{\linewidth}{|>{\raggedright\arraybackslash}m{4cm}|X|}
    \hline
    \textbf{Test ID} & utest\_2 \\
    \hline
    \textbf{Expected result} &
    \textbf{Test loosing heart via getting question wrong}
    \begin{itemize}
        \item Player loses 1 heart.
        \item Player’s health bar has a visual indicator of losing a single heart.
    \end{itemize}
    \textbf{Test loosing heart via getting hit by enemy}
    \begin{itemize}
        \item Player loses 1 heart.
        \item Player’s health bar has a visual indicator of losing a single heart.
    \end{itemize}
    \textbf{Test regaining heart via item}
    \begin{itemize}
        \item Player gains 1 heart.
        \item Heart recovered does not go over max hearts.
        \item Player’s health bar has a visual indicator of gaining a heart.
    \end{itemize}
    \\
    \hline
    \textbf{Pass/Fail} & Pass \\
    \hline
    \textbf{Severity (If Fail)} & Moderate: Failing this may lead to disruption to gameplay balance. However, this will not disrupt other essential features of the game. \\
    \hline
    \textbf{Comments} & None \\
    \hline
\end{tabularx}
\end{center}

\subsubsection{Collision Detection}

\begin{center}
\begin{tabularx}{\linewidth}{|>{\raggedright\arraybackslash}m{4cm}|X|}
    \hline
    \textbf{Test ID} & utest\_3 \\
    \hline
    \textbf{Description of test} & This test case verifies the collision detection between different objects, entities and maps in the game. The player should be restricted from moving outside the map boundaries, and collision objects should intuitively align with the player’s expectations based on natural in-game logic. \\
    \hline
    \textbf{Related requirement document details} & Movement Mechanics 2.2.8 \\
    \hline
    \textbf{Prerequisites for test} & None \\
    \hline
    \textbf{Test Procedure} &
    \textbf{Test player collisions with static objects}
    \begin{itemize}
        \item Have static objects present in the game world.
        \item Move the player character towards the object.
    \end{itemize}
    \textbf{Test player collision with dynamic objects}
    \begin{itemize}
        \item Have dynamic objects present in the game world.
        \item Move the player character towards the object.
    \end{itemize}
    \textbf{Test player movement near map boundaries}
    \begin{itemize}
        \item Have player characters and game maps loaded.
        \item Move the player character towards boundaries (e.g walls) of the map using movement system.
    \end{itemize}
    \\
    \hline
    \textbf{Test material used} & Keyboard, Xbox Controller \\
    \hline
\end{tabularx}
\end{center}

\begin{center}
\begin{tabularx}{\linewidth}{|>{\raggedright\arraybackslash}m{4cm}|X|}
    \hline
    \textbf{Expected result} &
    \textbf{Test player collisions with static objects}
    \begin{itemize}
        \item Player character should stop moving upon contact with the static object.
        \item Player character should be unable to pass through it.
        \item There should be no clipping or glitching of the character model.
    \end{itemize}
    \textbf{Test player collision with dynamic objects}
    \begin{itemize}
        \item Player character should either interact with or be blocked by the dynamic object, depending on the object’s movement or state (e.g standing on a moving platform).
        \item The collision response should appear natural, without any erratic movements.
    \end{itemize}
    \textbf{Test player movement near map boundaries}
    \begin{itemize}
        \item Player character should not be able to move beyond defined map boundaries.
        \item The character should stop at the boundary without any visual error, no part of the player model should appear outside the map area.
    \end{itemize}
    \\
    \hline
    \textbf{Comments} & None \\
    \hline
    \textbf{Created by} & SL \\
    \hline
    \textbf{Test environment(s)} & Windows 11, Mac OS 15.3, Unity 2022.3.50f1 \\
    \hline
\end{tabularx}
\end{center}

\noindent \textbf{Results:}

\begin{center}
\begin{tabularx}{\linewidth}{|>{\raggedright\arraybackslash}m{4cm}|X|}
    \hline
    \textbf{Test ID} & utest\_3 \\
    \hline
    \textbf{Expected result} &
    \textbf{Test player collisions with static objects}
    \begin{itemize}
        \item Player character should stop moving upon contact with the static object.
        \item Player character should be unable to pass through it.
        \item There should be no clipping or glitching of the character model.
    \end{itemize}
    \textbf{Test player collision with dynamic objects}
    \begin{itemize}
        \item Player character should either interact with or be blocked by the dynamic object, depending on the object’s movement or state.
        \item The collision response should appear natural, without any erratic movements.
    \end{itemize}
    \textbf{Test player movement near map boundaries}
    \begin{itemize}
        \item Player character should not be able to move beyond defined map boundaries.
        \item The character should stop at the boundary without any visual error, no part of the player model should appear outside the map area.
    \end{itemize}
    \\
    \hline
\end{tabularx}
\end{center}

\newpage

\begin{center}
\begin{tabularx}{\linewidth}{|>{\raggedright\arraybackslash}m{4cm}|X|}
    \hline
    \textbf{Regression 1 Result} &
    \textbf{Test player collisions with static objects}
    \begin{itemize}
        \item Player character immediately stops upon contact with static objects.
        \item Player cannot pass through the object under any condition.
        \item No clipping, glitching, or overlapping of the player model with the static object occurs.
    \end{itemize}
    \textbf{Test player collision with dynamic objects}
    \begin{itemize}
        \item Player either interacts with or is blocked by the dynamic object.
        \item Collision response is consistent with game mechanics.
    \end{itemize}
    \textbf{Test player movement near map boundaries}
    \begin{itemize}
        \item Player is blocked from moving beyond the map boundaries.
        \item No visual errors occur at the boundaries, player does not clip outside the map.
    \end{itemize}
    \textbf{Unexpected Results}
    \begin{itemize}
        \item No unexpected results
    \end{itemize}
    \\
    \hline
    \textbf{Pass/Fail} & Pass \\
    \hline
    \textbf{Severity (If Fail)} & High: Failure would disrupt core movement mechanics and the player may exploit bugs to bypass challenges or become stuck outside of the map. \\
    \hline
    \textbf{Comments} & None \\
    \hline
\end{tabularx}
\end{center}

\subsubsection{Enemy AI Behaviour}

\begin{center}
\begin{tabularx}{\linewidth}{|>{\raggedright\arraybackslash}m{4cm}|X|}
    \hline
    \textbf{Test ID} & utest\_4 \\
    \hline
    \textbf{Description of test} & This test case verifies the behavior of the enemy AI in a game environment, specifically focusing on their patrol and detection routines. The AI should follow a set patrol path and detect the player within a specified detection radius. \\
    \hline
    \textbf{Related requirement document details} & \makecell[l] {
        AI Boss (2.2.4) \\
        Combat Mechanics (2.2.9)
    } \\
    \hline
    \textbf{Prerequisites for test} & \makecell[l] {
        Enemy prefab created with AI controller. \\
        Test environment with patrol routes and obstacles to test pathfinding.
    } \\
    \hline
    \textbf{Test procedure} &
    \textbf{Test enemy patrol behavior}
    \begin{itemize}
        \item Place the enemy in the environment with a predefined patrol route.
        \item Observe the enemy moving smoothly between patrol points while scanning for the player within its field of view.
        \item Ensure the enemy turns around or adjusts its patrol when it reaches the end of the route or encounters a barrier.
    \end{itemize}
    \textbf{Test enemy detection of the player}
    \begin{itemize}
        \item Position the player within the enemy’s field of view.
        \item Verify the enemy enters a detection state when the player is visible.
        \item Ensure the detection is consistent with the game’s design (e.g., player being fully visible vs. partially obscured).
    \end{itemize}
    \textbf{Test edge cases for player visibility}
    \begin{itemize}
        \item Position the player at partially obscured locations (e.g., just outside the enemy’s field of view or behind thin obstacles).
        \item Verify that the enemy does not detect the player unless they are within the intended visibility range or in line of sight.
        \item Confirm that if the player moves quickly through the enemy’s field of view, the enemy reacts briefly before resuming its patrol.
    \end{itemize}
    \textbf{Test reaction to unreachable players}
    \begin{itemize}
        \item Place the player in positions the enemy cannot reach (e.g., higher platforms or behind barriers).
        \item Ensure the enemy does not attempt impossible paths to reach the player.
        \item Verify the enemy either paces in frustration, searches nearby, or returns to its patrol route after a delay.
    \end{itemize}
    \\
    \hline
    \textbf{Test material used} & Unity Console Logs, Unity Gizmos to see Raycasts \\
    \hline
\end{tabularx}
\end{center}

\begin{center}
\begin{tabularx}{\linewidth}{|>{\raggedright\arraybackslash}m{4cm}|X|}
    \hline
    \textbf{Expected result} &
    \textbf{Enemy Patrol Behavior:}
    \begin{itemize}
        \item The enemy moves smoothly between patrol points without erratic or unexpected movements.
        \item When reaching the end of the patrol route or encountering a barrier, the enemy turns around or adjusts direction appropriately.
        \item The enemy continuously scans for the player within its field of view.
    \end{itemize}
    \textbf{Enemy Detection of the Player:}
    \begin{itemize}
        \item The enemy immediately enters a detection state when the player enters its field of view.
        \item The detection should be consistent and only occur if the player is fully visible or within the specified visibility range.
    \end{itemize}
    \textbf{Enemy Behavior When the Player Hides:}
    \begin{itemize}
        \item If the player enters a designated stealth area, the enemy does not detect them, even if the enemy is nearby.
        \item The enemy investigates the player’s last known position before resuming patrol or searching elsewhere.
    \end{itemize}
    \textbf{Edge Cases for Player Visibility:}
    \begin{itemize}
        \item The enemy does not detect the player if they are outside its field of view or partially obscured as per game design rules.
        \item If the player moves quickly through the enemy’s field of view, the enemy reacts briefly but returns to patrol mode if the player disappears from sight.
    \end{itemize}
    \textbf{Reaction to Unreachable Players:}
    \begin{itemize}
        \item The enemy does not attempt impossible paths to reach the player in unreachable locations (e.g., on higher platforms or behind barriers).
        \item The enemy exhibits appropriate alternative behavior, such as pacing, investigating nearby areas, or returning to its patrol route after a short delay.
    \end{itemize}
    \\
    \hline
    \textbf{Created by} & WM \\
    \hline
    \textbf{Test environment(s)} & Windows 11, Mac OS 15.3, Unity 2022.3.50f1 \\
    \hline
\end{tabularx}
\end{center}

\newpage

\noindent \textbf{Results:}

\begin{center}
\begin{tabularx}{\linewidth}{|>{\raggedright\arraybackslash}m{4cm}|X|}
    \hline
    \textbf{Test ID} & utest\_4 \\
    \hline
    \textbf{Expected result} &
    \textbf{Enemy Patrol Behavior:}
    \begin{itemize}
        \item The enemy moves smoothly between patrol points without erratic or unexpected movements.
        \item When reaching the end of the patrol route or encountering a barrier, the enemy turns around or adjusts direction appropriately.
        \item The enemy continuously scans for the player within its field of view.
    \end{itemize}
    \textbf{Enemy Detection of the Player:}
    \begin{itemize}
        \item The enemy immediately enters a detection state when the player enters its field of view.
        \item The detection should be consistent and only occur if the player is fully visible or within the specified visibility range.
    \end{itemize}
    \textbf{Enemy Behavior When the Player Hides:}
    \begin{itemize}
        \item If the player enters a designated stealth area, the enemy does not detect them, even if the enemy is nearby.
        \item The enemy investigates the player’s last known position before resuming patrol or searching elsewhere.
    \end{itemize}
    \textbf{Edge Cases for Player Visibility:}
    \begin{itemize}
        \item The enemy does not detect the player if they are outside its field of view or partially obscured as per game design rules.
        \item If the player moves quickly through the enemy’s field of view, the enemy reacts briefly but returns to patrol mode if the player disappears from sight.
    \end{itemize}
    \textbf{Reaction to Unreachable Players:}
    \begin{itemize}
        \item The enemy does not attempt impossible paths to reach the player in unreachable locations (e.g., on higher platforms or behind barriers).
        \item The enemy exhibits appropriate alternative behavior, such as pacing, investigating nearby areas, or returning to its patrol route after a short delay.
    \end{itemize}
    \\
    \hline
    \textbf{Result} &
    \textbf{Enemy Patrol Behavior:}
    \begin{itemize}
        \item The enemy moves smoothly between patrol points without erratic or unexpected movements.
        \item When reaching the end of the patrol route or encountering a barrier, the enemy turns around or adjusts its direction appropriately.
        \item The enemy continuously scans for the player within its field of view as shown in the console logs.
    \end{itemize}
    \\
    \hline
\end{tabularx}
\end{center}

\begin{center}
\begin{tabularx}{\linewidth}{|>{\raggedright\arraybackslash}m{4cm}|X|}
    \hline
    \textbf{Result} &
    \textbf{Enemy Detection of the Player:}
    \begin{itemize}
        \item The enemy immediately enters a detection state when the player enters its field of view.
        \item Detection occurs only if the player is fully visible or within the specified visibility range, with no inconsistencies in detection.
    \end{itemize}
    \textbf{Enemy Behavior When the Player Hides:}
    \begin{itemize}
        \item If the player enters a designated stealth area, the enemy does not detect them, even if the enemy is nearby.
    \end{itemize}
    \textbf{Edge Cases for Player Visibility:}
    \begin{itemize}
        \item The enemy does not detect the player if they are outside its field of view or partially obscured.
    \end{itemize}
    \textbf{Reaction to Unreachable Players:}
    \begin{itemize}
        \item The enemy does not attempt impossible paths to reach the player in unreachable locations (e.g., on higher platforms or behind barriers).
    \end{itemize}
    \\
    \hline
    \textbf{Pass/Fail} & Partial Pass \\
    \hline
    \textbf{Severity (If Fail)} & Moderate \\
    \hline
    \textbf{Comments} & Partial player visibility detection had a few inconsistencies, where the enemy sometimes detected players behind thin obstacles unexpectedly. Further tuning required for AI visibility calculations. \\
    \hline
\end{tabularx}
\end{center}

\newpage

\subsubsection{Level Loading / Transitions}

\newpage

\subsubsection{SkillsBuild Questions}

\begin{center}
\begin{tabularx}{\linewidth}{|>{\raggedright\arraybackslash}m{4cm}|X|}
    \hline
    \textbf{Test ID} & utest\_6 \\
    \hline
    \textbf{Description of test} & Test that questions display and function correctly during gameplay. \\
    \hline
    \textbf{Related requirement document details} & Educations Engagement (2.2.2) \\
    \hline
    \textbf{Prerequisites for test} & Skills Build question bank is populated with questions and answers for each course. \\
    \hline
    \textbf{Test procedure} &
    \textbf{Display}
    \begin{itemize}
        \item Check that the quiz UI opens when a quiz is initiated and closes when the quiz is complete.
        \item Verify that the questions and answers are rendered correctly.
    \end{itemize}
    \textbf{Function}
    \begin{itemize}
        \item Ensure that when an answer is clicked the game checks if the answer is correct or incorrect and notifies the user accordingly.
        \item Assert that the user can only move on to the next question after answering the current question.
    \end{itemize}
    \\
    \hline
    \textbf{Test material used} & Skills Build question bank \\
    \hline
    \textbf{Expected result} &
    \textbf{Display}
    \begin{itemize}
        \item The quiz UI opens and closes appropriately.
        \item Questions and answers are rendered correctly.
    \end{itemize}
    \textbf{Function}
    \begin{itemize}
        \item Appropriate feedback is provided after answering a question.
        \item User progresses after answering a question.
    \end{itemize}
    \\
    \hline
    \textbf{Comments} & Integration, engagement and correctness of quizzes are assessed in uattest\_5. \\
    \hline
    \textbf{Created by} & BM \\
    \hline
    \textbf{Test environment(s)} & Windows 11, Mac OS 15.3, Unity 2022.3.50f1 \\
    \hline
\end{tabularx}
\end{center}

\newpage

\noindent \textbf{Results:}

\begin{center}
\begin{tabularx}{\linewidth}{|>{\raggedright\arraybackslash}m{4cm}|X|}
    \hline
    \textbf{Test ID} & utest\_6 \\
    \hline
    \textbf{Result} &
    \textbf{Display}
    \begin{itemize}
        \item The quiz UI opens and closes as expected.
        \item Questions and answers are rendered properly.
    \end{itemize}
    \textbf{Function}
    \begin{itemize}
        \item The correct choice of feedback is displayed to the user.
        \item It is only possible to progress to the next question after answering it.
    \end{itemize}
    \\
    \hline
    \textbf{Pass/Fail} & Pass \\
    \hline
    \textbf{Severity (If Fail)} & Critical: If for some reason a question has no correct answers the player cannot progress thus game progression is impossible. \\
    \hline
    \textbf{Comments} & None \\
    \hline
\end{tabularx}
\end{center}

\newpage

\subsection{System Testing}

\subsubsection{Graphics / Audio}

\begin{center}
\begin{tabularx}{\linewidth}{|>{\raggedright\arraybackslash}m{4cm}|X|}
    \hline
    \textbf{Test ID} & stest\_1 \\
    \hline
    \textbf{Description of test} & Graphics and Audio: Verify that graphics render smoothly, textures load correctly, and audio plays without distortion or lag \\
    \hline
    \textbf{Related requirement document details} & \makecell[l]{
        SFX (2.2.5) \\
        Stylised Cold War Graphics (2.2.6)
    } \\
    \hline
    \textbf{Prerequisites for test} & Game is installed, graphics and audio are (partially) implemented, audio / display devices available, performance monitoring tools. \\
    \hline
    \textbf{Test procedure} &
    \textbf{Graphics Performance:}
    \begin{itemize}
        \item Verify consistent performance during complex scene transitions
        \item Verify rendering at multiple resolutions
    \end{itemize}
    \textbf{Audio systems performance}
    \begin{itemize}
        \item Check sound effects synchronize with animations
        \item Verify audio quality during transitions
    \end{itemize}
    \\
    \hline
    \textbf{Test material used} & \makecell[l]{
        Multiple devices \\
        Audio monitoring tools \\
        Performance profiler
    }
    \\
    \hline
    \textbf{Expected result} &
    \textbf{Graphics performance}
    \begin{itemize}
        \item Smooth transitions
        \item Consistent rendering across platforms
    \end{itemize}
    \textbf{Audio systems}
    \begin{itemize}
        \item Audio and sounds effects in sync
        \item Audio quality remains consistent
    \end{itemize}
    \\
    \hline
    \textbf{Comments} & None \\
    \hline
    \textbf{Created By} & DM \\
    \hline
    \textbf{Test environment(s)} & Windows 11, Mac OS 15.3, Unity 2022.3.50f1 \\
    \hline
\end{tabularx}
\end{center}

\newpage

\subsubsection{UI Navigation (Saving / Loading)}

\begin{center}
\begin{tabularx}{\linewidth}{|>{\raggedright\arraybackslash}m{4cm}|X|}
    \hline
    \textbf{Test ID} & stest\_2 \\
    \hline
    \textbf{Description of test} & UI Navigation (Saving and Loading): Test that the UI allows seamless navigation, and game states save/load correctly \\
    \hline
    \textbf{Related requirement document details} & Pause Menu (2.2.3) \\
    \hline
    \textbf{Prerequisites for test} & \makecell[l]{
        Complete UI implementation \\
        Save/Load system implemented
    }
    \\
    \hline
    \textbf{Test procedure} &
    \textbf{UI Navigation Flow:}
    \begin{itemize}
        \item Check navigation works correctly between all menu screens (main menu, options, save/load screens)
        \item Test menu responsiveness and input handling across different input methods
    \end{itemize}
    \textbf{Save system validation:}
    \begin{itemize}
        \item Test automatic save functionality throughout different parts of the game
        \item Verify manual save operations across different parts
        \item Validate correct restoration of game progress
    \end{itemize}
    \textbf{Load System Verification:}
    \begin{itemize}
        \item Test loading saves from different saves
        \item Verify handling of corrupted or incompatible save files
    \end{itemize}
    \\
    \hline
    \textbf{Test material used} & \makecell[l] {
        UI interaction logs \\
        Xbox controller
    } \\
    \hline
    \textbf{Expected result} &
    \begin{itemize}
        \item Smooth and intuitive UI navigation
        \item Reliable save file creation and management
        \item Accurate game state restoration when loading
    \end{itemize}
    \\
    \hline
    \textbf{Comments} & None \\
    \hline
    \textbf{Created by} & DM \\
    \hline
    \textbf{Test environment(s)} & Windows 11, Mac OS 15.3, Unity 2022.3.50f1 \\
    \hline
\end{tabularx}
\end{center}

\newpage

\subsubsection{Gameplay Flow}

\newpage

\subsubsection{Character Interactions}

\begin{center}
\begin{tabularx}{\linewidth}{|>{\raggedright\arraybackslash}m{4cm}|X|}
    \hline
    \textbf{Test ID} & stest\_4 \\
    \hline
    \textbf{Description of test} & Test the interaction between the game character and NPC, including the dialog system, sneaking mechanism, and NPC's reaction to the player.Making it fits the game's story and integrates smoothly into IBM Skills Build-related content. \\
    \hline
    \textbf{Related requirement document details} & \makecell[l]{
        Educational Engagement (2.2.2) \\
        Mission System (2.2.7)
    } \\
    \hline
    \textbf{Prerequisites for test} & There needs to be at least a preliminary complete question bank for teaching IBM knowledge and a simplified version of the game with full character interaction mechanics for testing basic functionality \\
    \hline
    \textbf{Test procedure} &
    \begin{itemize}
        \item Start the game and enter the section that contains the NPC and try the mission objectives.
        \item Engage in dialogue with the NPC and observe whether the IBM Skills Build questions are triggered naturally.
        \item Observe whether the character interactions including accepting tasks, interacting with the environment, communicating with enemies, check if they run smoothly and affect the progress of the game normally.
    \end{itemize}
    \\
    \hline
    \textbf{Test material used} & Laptop \\
    \hline
    \textbf{Expected result} &
    \begin{itemize}
        \item NPC interaction text is on-topic and does not introduce learning content out of the blue.
        \item Players are able to effectively access IBM-related issues through character interactions and can influence the progress of the game.
        \item The character's basic movement commands can cause interactions with NPCs, such as sneak and be found, also the dialogue with NPC works normally.
    \end{itemize}
    \\
    \hline
    \textbf{Comments} & None \\
    \hline
    \textbf{Created by} & HC \\
    \hline
    \textbf{Test environment(s)} & Windows 11, Mac OS 15.3, Unity 2022.3.50f1 \\
    \hline
\end{tabularx}
\end{center}

\newpage

\subsection{User Acceptance Testing}

\subsubsection{Gameplay Experience}

\begin{center}
\begin{tabularx}{\linewidth}{|>{\raggedright\arraybackslash}m{4cm}|X|}
    \hline
    \textbf{Test ID} & uattest\_1 \\
    \hline
    \textbf{Description of test} & Testing the overall game experience, including sneaking mechanics, combat system, mission flow,AI, and the way IBM Skills Build content is embedded to ensure the balance of education and entertainment. \\
    \hline
    \textbf{Related requirement document details} & \makecell[l]{
        Stealth Gameplay (2.2.1) \\
        Combat Mechanics (2.2.9) \\
        Educational Engagement (2.2.2)
    } \\
    \hline
    \textbf{Prerequisites for test} & Executable version of the game (with sneaking, combat, quiz mechanics) with complete level design with missions, enemies, IBM-related challenges and a simple AI \\
    \hline
    \textbf{Test procedure} &
    \begin{itemize}
        \item Enter a level that includes sneaking and fighting.
        \item Test the smoothness of the game operation to ensure that there is no delay or abnormality in the switching between sneaking, fighting and answering questions.
        \item Observe how IBM Skills Build related questions are triggered to make sure that they do not affect the player's immersion or increase the difficulty of the game too much.
    \end{itemize}
    \\
    \hline
    \textbf{Test material used} & Laptop \\
    \hline
    \textbf{Expected result} &
    \begin{itemize}
        \item The game is smooth and can run normally.
        \item Players are able to sneak and fight without any problems, and the learning content does not spoil the game experience.
        \item The game's upgraded mechanics will keep the player's interest for educational purposes.
        \item The AI is able to adjust the difficulty of the game.
    \end{itemize}
    \\
    \hline
    \textbf{Comments} & None \\
    \hline
    \textbf{Created by} & HC \\
    \hline
    \textbf{Test environment(s)} & Windows 11, Mac OS 15.3, Unity 2022.3.50f1 \\
    \hline
\end{tabularx}
\end{center}

\newpage

\subsubsection{User Interface}

\begin{center}
\begin{tabularx}{\linewidth}{|>{\raggedright\arraybackslash}m{4cm}|X|}
    \hline
    \textbf{Test ID} & uattest\_2 \\
    \hline
    \textbf{Description of test} & Test the game UI to see if it is intuitive, if the information display is clear, if the quiz interface is in line with the overall style of the game, and to test the stability of the UI to ensure that each function works successfully. Ensure that players can use and experience IBM Skills Build content smoothly. \\
    \hline
    \textbf{Related requirement document details} & \makecell[l]{
        Pause Screen (2.2.3) \\
        Stylised Cold War Graphics (2.2.6)
    } \\
    \hline
    \textbf{Prerequisites for test} & Complete UI design and interaction logic has been implemented. In-game question and answer interface, mission menu, pause menu and other functions are available. \\
    \hline
    \textbf{Test procedure} &
    \begin{itemize}
        \item Enter the main interface of the game, check if the navigation menu is clear and the main functions are easy to find.
        \item Enter the mission level and check whether the UI feedback (such as mission guidelines, blood level, props information) is intuitive.
        \item Enter the IBM Skills Build question interface and check whether the interaction is smooth and the UI design is in line with the overall style of the game.
        \item Repeat testing the buttons in the UI to test their stability and whether they can interact with each other properly.
    \end{itemize}
    \\
    \hline
    \textbf{Test material used} & Laptop \\
    \hline
    \textbf{Expected result} &
    \begin{itemize}
        \item The main menu of the game, the mission interface and the answer interface have the same style, and the operation is intuitive.
        \item UI feedback is timely, important information will not be obscured by other interface elements.
        \item The function of buttons in the UI runs normally and stably, and can interact with other buttons.
    \end{itemize}
    \\
    \hline
    \textbf{Comments} & None \\
    \hline
    \textbf{Created by} & HC \\
    \hline
    \textbf{Test environment(s)} & Windows 11, Mac OS 15.3, Unity 2022.3.50f1 \\
    \hline
\end{tabularx}
\end{center}

\newpage

\subsubsection{Performance and Stability}

\begin{center}
\begin{tabularx}{\linewidth}{|>{\raggedright\arraybackslash}m{4cm}|X|}
    \hline
    \textbf{Test ID} & uattest\_3 \\
    \hline
    \textbf{Description of test} &
    Application must perform well on low-end hardware, must not exhibit sudden framedrops when performing and can run indefinitely without hanging, crashing or leaking memory.
    \\
    \hline
    \textbf{Related requirement document details} & \makecell[l]{
        Risks and Issues (3.1)
    } \\
    \hline
    \textbf{Prerequisites for test} &
        Build of the game with majority of the features implemented.
    \\
    \hline
    \textbf{Test procedure} &
    \textbf{Performance}
    \begin{itemize}
        \item Play the game for a while and monitor the average frame rate.
        \item Check the latency graph for framedrops.
    \end{itemize}
    \textbf{Stability}
    \begin{itemize}
        \item Play the game for an extended period of time, change levels, settings, difficulty.
        \item Leave the game idle for an extended period of time, ensure CPU, GPU and memory usage do not increase over time.
    \end{itemize}
    \\
    \hline
    \textbf{Test material used} & Laptop (with integrated graphics) \\
    \hline
    \textbf{Expected result} &
    \textbf{Performance}
    \begin{itemize}
        \item Game runs at an acceptable framerate ($60 \geq$), with minimal changes in frame rate when performing actions.
        \item There are no perceptible framerate drops
    \end{itemize}
    \textbf{Stability}
    \begin{itemize}
        \item System resource usage remains consistent
    \end{itemize}
    \\
    \hline
    \textbf{Comments} & \makecell[l]{
        It may require more advanced graphical pipeline debugging to identify \\ framedrops. \\ \\
        Stability testing is time consuming and realistically, this test will not be \\ run for more than 30 minutes at a time.
    } \\
    \hline
    \textbf{Created by} & TD \\
    \hline
    \textbf{Test environment(s)} & Windows 11, Mac OS 15.3, Unity 2022.3.50f1 \\
    \hline
\end{tabularx}
\end{center}

\newpage

\subsubsection{Accessibility}

\begin{center}
\begin{tabularx}{\linewidth}{|>{\raggedright\arraybackslash}m{4cm}|X|}
    \hline
    \textbf{Test ID} & uattest\_4 \\
    \hline
    \textbf{Description of test} & This test case verifies the game's accessibility features, focusing on customizable controls, subtitles, colorblind-friendly options, and audio cues for visually impaired players. \\
    \hline
    \textbf{Related requirement document details} & N/A \\
    \hline
    \textbf{Prerequisites for test} & Accessibility features implemented, game settings menu fully functional \\
    \hline
    \textbf{Test procedure} &
    \textbf{Customizable controls:}
    \begin{itemize}
        \item Check that players can remap all input keys and save their preferences.
    \end{itemize}
    \textbf{Subtitles}
    \begin{itemize}
        \item Verify that subtitles are available, accurately synced with dialogues, and customizable in size and color.
    \end{itemize}
    \textbf{Colorblind mode:}
    \begin{itemize}
        \item Enable colorblind-friendly mode and confirm that game elements are distinguishable for various types of colorblindness.
    \end{itemize}
    \textbf{Audio Cues:}
    \begin{itemize}
        \item Test that all critical game events have corresponding audio cues.
    \end{itemize}
    \\
    \hline
    \textbf{Test material used} & Accessibility test checklist, user input logs \\
    \hline
    \textbf{Expected result} & Players can fully customize controls and subtitle preferences. Colorblind mode offers distinguishable visuals. Audio cues align with game events to assist visually impaired players. \\
    \hline
    \textbf{Comments} & None \\
    \hline
    \textbf{Created by} & WM \\
    \hline
    \textbf{Test environment(s)} & Windows 11, Mac OS 15.3, Unity 2022.3.50f1 \\
    \hline
\end{tabularx}
\end{center}

\newpage

\noindent \textbf{Results:}

\begin{center}
\begin{tabularx}{\linewidth}{|>{\raggedright\arraybackslash}m{4cm}|X|}
    \hline
    \textbf{Test ID} & uattest\_4 \\
    \hline
    \textbf{Expected result} & Players can fully customize controls and subtitle preferences. Colorblind mode offers distinguishable visuals. Audio cues align with game events to assist visually impaired players. \\
    \hline
    \textbf{Result} &
    \textbf{Customizable Controls:}
    \begin{itemize}
        \item Control Scheme changed and functioned as expected
    \end{itemize}
    \textbf{Subtitles:}
    \begin{itemize}
        \item Cover most elements, sometimes slightly out of time
    \end{itemize}
    \textbf{Colorblind Mode:}
    \begin{itemize}
        \item Provides distinguishable elements for most cases but has issues with red-green differentiation
    \end{itemize}
    \textbf{Audio Cues:}
    \begin{itemize}
        \item Cover most elements, sometimes slightly out of time and a few minor scenarios.
    \end{itemize}
    \\
    \hline
    \textbf{Pass/Fail} & Partial Pass \\
    \hline
    \textbf{Severity (If Fail)} & Moderate \\
    \hline
    \textbf{Comments} & Additional testing and adjustments needed for red-green colorblind compatibility, subtitle timing, and missing audio cues during boss encounters. \\
    \hline
    \textbf{Test environment(s)} & Windows 11, Mac OS 15.3, Unity 2022.3.50f1 \\
    \hline
\end{tabularx}
\end{center}

\newpage

\subsubsection{Quizzes}

\begin{center}
\begin{tabularx}{\linewidth}{|>{\raggedright\arraybackslash}m{4cm}|X|}
    \hline
    \textbf{Test ID} & uattest\_5 \\
    \hline
    \textbf{Description of test} & Ensure the quizzes are engaging, well-integrated and correctly assessed. \\
    \hline
    \textbf{Related requirement document details} & Education Engagement (2.2.2) \\
    \hline
    \textbf{Prerequisites for test} & A functional game that passes utest\_2, utest\_6, utest\_10, utest\_13 and utest\_14. \\
    \hline
    \textbf{Test procedure} & Load the game and select a course. Play through the game until a quiz is initiated.

    \textbf{Engaging:}
    \begin{itemize}
        \item Verify that both audio and visual feedback are provided after correct and incorrect answers.
    \end{itemize}
    \textbf{Well-integrated:}
    \begin{itemize}
        \item Ensure a smooth transition from gameplay to the quiz and vice versa.
        \item Check that the quiz UI matches the game design and theme.
    \end{itemize}
    \textbf{Correctly assessed:}
    \begin{itemize}
        \item Check that the questions relate to the course selected.
        \item Answer a question incorrectly and check that the correct answer is shown.
    \end{itemize}
    \\
    \hline
    \textbf{Test material used} & Skills Build question bank \\
    \hline
    \textbf{Expected result} & 
    \textbf{Engaging:}
    \begin{itemize}
        \item Audio and visual feedback is provided after answering correctly or incorrectly.
    \end{itemize}
    \textbf{Well-integrated:}
    \begin{itemize}
        \item There is a smooth transition from gameplay to the quiz and vice versa.
        \item The quiz UI matches the game design and theme.
    \end{itemize}
    \textbf{Correctly assessed:}
    \begin{itemize}
        \item The questions relate to the course selected.
        \item The correct answer is shown when the user answers incorrectly.
    \end{itemize} 
    \\
    \hline
    \textbf{Comments} & Ensure edge cases such as a player losing all their health during a quiz are handled. \\
    \hline
    \textbf{Created by} & BM \\
    \hline
    \textbf{Test environment(s)} & Windows 11, Mac OS 15.3, Unity 2022.3.50f1 \\
    \hline
\end{tabularx}
\end{center}

\noindent \textbf{Results:}

\begin{center}
\begin{tabularx}{\linewidth}{|>{\raggedright\arraybackslash}m{4cm}|X|}
    \hline
    \textbf{Test ID} & uattest\_5 \\
    \hline
    \textbf{Result} &
    \textbf{Engaging:}
    \begin{itemize}
        \item Audio and visual feedback is provided after answering correctly or incorrectly making the quizzes more engaging.
    \end{itemize}
    \textbf{Well-integrated:}
    \begin{itemize}
        \item There is a smooth transition from gameplay to the quiz and from the quiz to gameplay.
        \item The quiz UI matches the game design and theme.
    \end{itemize}
    \textbf{Correctly assessed:}
    \begin{itemize}
        \item The questions relate to the course selected.
        \item The correct answer is shown when the user answers incorrectly, helping them learn from their mistakes.
    \end{itemize} 
    \\
    \hline
    \textbf{Pass/Fail} & Pass \\
    \hline
    \textbf{Severity (If Fail)} & Moderate for engagement and integration but High for correctly assessed as it would affect core gameplay if it fails. \\
    \hline
    \textbf{Comments} & None. \\
    \hline
\end{tabularx}
\end{center}

\newpage

\section{Testing Context}

\noindent To ensure compatibility and coverage across educational environments, the following tools and environments are used for testing the system:

\subsection{Operating Systems}

\begin{itemize}
    \item Windows 11: Selected due to its widespread use in educational institutions and professional environments. Ensures compatibility for schools and educational labs.
    \item macOS (15.3 and above): Included as many schools and creative institutions prefer macOS systems for educational and development purposes.
\end{itemize}

\subsection{Testing Tools and Drivers}

\begin{itemize}
    \item Unity Editor (2022.3.50f1): Used to test and debug game mechanics, AI behaviors, and user interactions.
    \item Accessibility Testing Tools: Simulators for colorblindness and screen reader compatibility (e.g., Color Oracle, VoiceOver for macOS).
    \item Input Device Testing: Keyboard and controller compatibility tests using standard USB and wireless game controllers.
\end{itemize}

\subsection{Test Failure Severity Levels}

\begin{center}
\begin{tabularx}{\linewidth}{|c|X|}
    \hline
    \textbf{Severity Level} & \textbf{Description} \\
    \hline
    \textbf{Critical} & The system crashes or produces incorrect outcomes that prevent gameplay or educational content delivery. Requires immediate attention. \\
    \hline
    \textbf{High} & Major functionality is disrupted, such as inaccessible accessibility features or broken AI behaviors, affecting core gameplay. \\
    \hline
    \textbf{Moderate} & Gameplay is impacted but still functional, such as minor UI glitches, missing audio cues, or partial AI inconsistencies. \\
    \hline
    \textbf{Low} & Minor inconveniences that do not significantly impact the user experience, such as aesthetic inconsistencies or delayed animations. \\
    \hline
\end{tabularx}
\end{center}

\end{document}
